\documentclass[10pt,landscape]{article}
\usepackage{multicol}
\usepackage{calc}
\usepackage{ifthen}
\usepackage[landscape, a4paper, top=3mm, left=3mm, right=3mm, bottom=3mm]{geometry}
\usepackage[
     pdftex,
     bookmarksopen=true,
     pdfauthor={Justin Iven Müller},
     pdftitle={Theoretische Informatik Cheatsheet},
     colorlinks=true,
     linkcolor=blue,
     urlcolor=blue
]{hyperref}
\usepackage{amsmath,amssymb,amsfonts}
\usepackage{enumitem}
\usepackage{tikz}
\usetikzlibrary{shapes,positioning,arrows,fit,calc,graphs,graphs.standard,intersections, automata}
\usepackage{pgfplots}
\usepackage[ngerman]{babel}
\usepackage{faktor}
\usepackage{stmaryrd}

\begin{document}
\input{content/def.tex}
\begin{multicols*}{4}

\begin{center}
     \Large{\textbf{Theoretische Informatik}} \\
\end{center}

\input{content/alphabetWordsLanguage.tex}
\input{content/regularExpressions.tex}
\section{Endliche Automaten}

\subsection{Grundidee}
Endliche Automaten (EA) sind speicherbeschränkte Berechnungsmodelle zur Erkennung formaler Sprachen.
Der einzige interne Speicher ist der aktuelle Zustand.
Ein Eingabewort wird sequenziell gelesen; Akzeptanz wird über Endzusände entschieden.

\subsection{Deterministische endliche Automaten}
Ein deterministischer endlicher Automat (DEA) ist ein Tupel
\[
    M = (Q, \Sigma, \delta, q_0, F)
\]
mit
\begin{itemize}
    \item endlicher Zustandsmenge \(Q\),
    \item Alphabet \(\Sigma\),
    \item totaler Übergangsfunktion \(\delta : Q \times \Sigma \to Q\),
    \item Startzustand \(q_0 \in Q\),
    \item Endzustandsmenge \(F \subseteq Q\).
\end{itemize}

\subsection{Konfigurationen}
Eine Konfiguration ist ein Paar
\[
    (q,w) \in Q \times \Sigma^*
\]
mit aktuellem Zustand \(q\) und Restwort \(w\).

\subsection{Berechnungsschritte}
Die Ein-Schritt-Relation ist definiert durch
\[
    (q,w) \vdash_M (p,x)
\]
genau dann, wenn \(w=ax\) mit \(a\in\Sigma\) und \(\delta(q,a)=p\).

Die reflexiv-transitive Hülle wird mit
\[
    (q,w) \vdash_M^* (p,x)
\]
bezeichnet.

\subsection{Akzeptanz und Sprache}
Für \(w\in\Sigma^*\) startet die Berechnung in \((q_0,w)\).
\(w\) wird akzeptiert, falls
\[
    (q_0,w) \vdash_M^* (q_f,\varepsilon) \quad\text{für ein } q_f\in F.
\]

Die erkannte Sprache ist
\[
    L(M)=\{\,w\in\Sigma^* \mid (q_0,w)\vdash_M^*(q_f,\varepsilon) \text{ für ein } q_f\in F\,\}.
\]

\subsection{Nichtdeterministische endliche Automaten}
Ein nichtdeterministischer endlicher Automat (NEA) hat dieselbe Struktur
\[
    M = (Q,\Sigma,\delta,q_0,F)
\]
wobei
\[
    \delta : Q \times \Sigma \to \mathcal{P}(Q)
\]
gilt. Damit sind mehrere Folgezustände pro Symbol möglich.

Akzeptanz: Es existiert mindestens ein Lauf, der in einem Zustand aus \(F\) endet.

\subsection{Äquivalenz von DEA und NEA}
Sein NEA \(N=(Q_N,\Sigma,\delta_N,q_{0N},F_N)\) kann in einen equivalenten DEA \(D=(Q_D,\Sigma,\delta_D,q_{0D},F_D)\) überführt werden durch
\begin{itemize}
    \item \(Q_D = \mathcal{P}(Q_N)\) (Mengen von NEA-Zuständen),
    \item \(F_D = \{S \subseteq Q_N \mid S \cap F_N \neq \emptyset\}\) (DEA-Endzustände sind Mengen, die mindestens einen NEA-Endzustand enthalten),
    \item \(\delta_D(S,a) = \bigcup_{q \in S} \delta_N(q,a)\) für \(S \subseteq Q_N\) und \(a \in \Sigma\) (Menge aller Zustände, die durch \(a\) von einem Zustand in \(S\) erreichbar sind),
\end{itemize}

\subsection{$\varepsilon$-Übergänge}
\(\varepsilon\)-NEAs erlauben Übergänge ohne Symbolkonsum mit
\[
    \delta : Q \times (\Sigma \cup \{\varepsilon\}) \to \mathcal{P}(Q).
\]
Auch \(\varepsilon\)-NEAs sind durch Eliminierung von \(\varepsilon\)-Übergängen und anschliessende Determinisierung zu DEA äquivalent.


\subsection{Äquivalenz DEA und RA}
Jede von einem DEA erkannte Sprache kann durch einen regulären Ausdruck beschrieben werden, und umgekehrt. Diese Sprachen werden als reguläre bezeichnet.
Reguläre Sprachen sind durch verschiedene Formalismen äquivalent beschreibbar, darunter DEA, NEA, \(\varepsilon\)-NEA und reguläre Ausdrücke.
\begin{itemize}
    \item \textbf{Akzeptierende Mech.:} DEA, NEA, \(\varepsilon\)-NEA
    \item \textbf{Beschreibungsform:} reguläre Ausdrücke
\end{itemize}

\subsection*{Abschlusseigenschaften regulärer Sprachen}
Die Klasse der regulären Sprachen ist unter verschiedenen Operationen abgeschlossen.

Seien $L, L_1, L_2 \subseteq \Sigma^*$ reguläre Sprachen. Dann sind ebenfalls regulär:
\begin{itemize}
    \item Vereinigung: \(L_1 \cup L_2\)
    \item Konkatenation: \(L_1 L_2\)
    \item Kleenesche Hülle: \(L^*\)
    \item Komplement: \(\overline{L} = \Sigma^* \setminus L\)
    \item Schnitt: \(L_1 \cap L_2\)
    \item Mengendifferenz: \(L_1 \setminus L_2\)
\end{itemize}

\subsection*{Zustandsklassen}
Zu jedem Zustand $p$ eines Automaten $M$ gehört die Zustandsklasse
\[
[p] = \{ w \in \Sigma^* \mid M \text{ endet nach dem Lesen von } w \text{ in } p \}.
\]

Die Zustandsklassen besitzen folgende Eigenschaften:
\begin{itemize}
    \item Die Menge aller Wörter wird durch die Klassen partitioniert:
    \[
    \Sigma^* = \bigcup_{p \in Q} [p]
    \]
    \item Für verschiedene Zustände sind die Klassen disjunkt:
    \[
    [p] \cap [q] = \emptyset \quad \text{für } p \ne q
    \]
    \item Die von $M$ akzeptierte Sprache ergibt sich als Vereinigung der Klassen akzeptierender Zustände:
    \[
    L(M) = \bigcup_{p \in F} [p]
    \]
\end{itemize}

\subsection*{Grenzen endlicher Automaten}
Ein endlicher Automat hat nur endlich viele Zustände und deshalb nur begrenzten Speicher.

Landen zwei Wörter $x$ und $y$ im selben Zustand, kann der Automat sie nicht mehr unterscheiden. Daher gilt für jedes Fortsetzungswort $z \in \Sigma^*$:
\[
xz \in L(M) \iff yz \in L(M).
\]

Wörter derselben Zustandsklasse verhalten sich also bei jeder Fortsetzung gleich.
Damit kann man Mindestzahlen für die benötigten Zustände eines Automaten ableiten.

\section{Kontextfreie Grammatiken}

\subsection{Definition}
Eine kontextfreie Grammatik (KFG) ist ein Tupel
\[
	G = (N, \Sigma, P, S)
\]
mit
\begin{itemize}
	\item Alphabet der \textbf{Nichtterminalen} $N$,
	\item Alphabet der \textbf{Terminalen} $\Sigma$ mit $N \cap \Sigma = \emptyset$,
	\item Endliche Menge \textbf{Produktionen} $P \subseteq N \times (N \cup \Sigma)^*$,
	\item \textbf{Startsymbol} $S \in N$.
\end{itemize}

Eine Produktion hat die Form
\[
	A \to \alpha \quad \text{mit } A \in N,\; \alpha \in (N \cup \Sigma)^*.
\]
\(\alpha \in (N \cup \Sigma)^*\) wird auch als \textbf{Satzformel} bezeichnet.

\subsection{Ableitung}
Der Ableitungsschritt ist definiert durch
\[
	uAv \Rightarrow_G u\alpha v
\]
genau dann, wenn $A \to \alpha \in P$ und $u,v \in (N \cup \Sigma)^*$.

Die reflexiv-transitive Hülle wird mit $\Rightarrow_G^*$ bezeichnet.

\subsection{Erzeugte Sprache}
Die von $G$ erzeugte Sprache ist
\[
	L(G) = \{\,w \in \Sigma^* \mid S \Rightarrow_G^* w\,\}.
\]
Eine Sprache $L \subseteq \Sigma^*$ heisst kontextfrei, falls es eine KFG $G$ mit $L=L(G)$ gibt.

\subsection{Links- und Rechtsableitung}
\begin{itemize}
	\item \textbf{Linksableitung:} In jedem Schritt wird das linkeste Nichtterminal ersetzt.
	\item \textbf{Rechtsableitung:} In jedem Schritt wird das rechteste Nichtterminal ersetzt.
\end{itemize}
Beide beschreiben dieselbe erzeugte Sprache.

\subsection{Ableitungsbaum}
Ein Ableitungsbaum repräsentiert die hierarchische Struktur einer Ableitung:
\begin{itemize}
	\item Wurzel ist $S$.
	\item Innere Knoten sind Nichtterminale.
	\item Blätter sind Terminalsymbole.
\end{itemize}
Die von links nach rechts gelesenen Blätter ergeben das erzeugte Wort.

\subsection{Mehrdeutigkeit}
Eine Grammatik ist mehrdeutig, wenn es ein Wort $w \in L(G)$ mit mindestens zwei verschiedenen Ableitungsbäumen gibt.
Eine Sprache heisst inherent mehrdeutig, wenn jede Grammatik, die sie erzeugt, mehrdeutig ist.

\section{Kellerautomaten}

\subsection{Grundidee}
Kellerautomaten (KA) erweitern endliche Automaten um einen Stack (Keller) als Speicherstruktur mit LIFO-Zugriff.
Sie erkennen genau die Klasse der kontextfreien Sprachen.

\subsection{Deterministischer Kellerautomat }
Ein deterministischer Kellerautomat ist ein Tupel
\[
    M = (Q, \Sigma, \Gamma, \delta, q_0, \$, F)
\]
wobei
\begin{itemize}
    \item \(Q\) eine endliche Menge von Zuständen,
    \item \(\Sigma\) ein endliches Eingabealphabet,
    \item \(\Gamma\) ein endliches Kelleralphabet,
    \item \(\delta : Q \times \Sigma \times \Gamma \to Q \times \Gamma^*\) eine Übergangsfunktion,
    \item \(q_0 \in Q\) der Startzustand,
    \item \(\$ \in \Gamma\) das Startkellersymbol,
    \item \(F \subseteq Q\) die Menge der akzeptierenden Zustände.
\end{itemize}

Damit der Determinismus gewährleistet ist, gilt für jeden Zustand \(q\) und alle Symbole \(a \in \Sigma\) und \(x \in \Gamma\), wenn \(\delta(q, a, x)\) definiert ist, dann ist \(\delta(q, \varepsilon, x)\) nicht definiert und umgekehrt.

\subsection{Graphische Darstellung}
Der Übergang \(\delta(q_0, a, x) = (q_1, w)\) wird grafisch dargestellt als\\
\begin{center}

    \begin{tikzpicture}[main/.style = {draw, circle}]
        \node[main] (q) {$q_0$};
        \node[main] (p) [right=1.5cm of q] {$q_1$};
        \draw[->] (q) to node[above] {\(a, x/w\)} (p);
    \end{tikzpicture}
\end{center}

\subsection{Nichtdeterministischer Kellerautomat}
Ein nichtdeterministischer Kellerautomat wird durch den selben Tupel wie der deterministische definiert, jedoch ist die Übergangsfunktion
\[
    \delta : Q \times (\Sigma \cup \{\varepsilon\}) \times \Gamma \to \mathcal{P}(Q \times \Gamma^*),
\]

\subsection{Konfigurationen}
Eine Konfiguration ist ein Tripel
\[
    (q,w,\gamma) \in Q \times \Sigma^* \times \Gamma^*
\]
mit aktuellem Zustand \(q\), Restwort \(w\) und Kellerinhalt \(\gamma\) (oberstes Symbol links).
\begin{itemize}
    \item \((q_0,w,\$)\) ist die Startkonfiguration.
    \item \((q,\varepsilon,\gamma)\) ist eine Endkonfiguration.
\end{itemize}

\subsection{Berechnungsschritt}
Ein Schritt
\[
    (q,aw,x\beta) \vdash_M (p,w,\gamma\beta)
\]
ist möglich, falls \((p,\gamma) \in \delta(q,a,x)\) für \(a \in \Sigma\).

Ein \(\varepsilon\)-Schritt
\[
    (q,w,x\beta) \vdash_M (p,w,\gamma\beta)
\]
ist möglich, falls \((p,\gamma) \in \delta(q,\varepsilon,x)\).

Die reflexiv-transitive Hülle wird mit \(\vdash_M^*\) bezeichnet.

\subsection{Akzeptanz}
Ein Wort \(w\) wird von \(M\) akzeptiert, wenn es eine Berechnung gibt, die von der Startkonfiguration \((q_0,w,\$)\) zu einer Endkonfiguration \((q_f,\varepsilon,\gamma)\) mit \(q_f \in F\) und \(\gamma \in \Gamma^*\) führt:
\[
    w \in L(M) \iff (q_0,w,\$) \vdash_M^* (q_f,\varepsilon,\gamma)
\]

\subsection{Beziehung zu Grammatiken}
Eine Sprache ist kontextfrei genau dann, wenn es einen nichtdeterministischen Kellerautomaten gibt, der sie akzeptiert.

\section{Turingmaschinen}

\subsection{Grundidee}
Eine Turing-Maschine (TM) ist ein endlicher Automat, ergänzt um einen Lese-/Schreibkopf und ein beidseitig unendliches Band von Zellen.
Die Eingabe steht zu Beginn auf dem Band; alle übrigen Zellen enthalten das Leerzeichen \(\shortsqcup\).

\subsection{Deterministische Turing-Maschine}
Eine deterministische TM ist ein 7-Tupel
\[
    M = (Q, \Sigma, \Gamma, \delta, q_0, \shortsqcup, F)
\]
mit
\begin{itemize}
    \item endlicher Zustandsmenge \(Q\),
    \item Eingabealphabet \(\Sigma\) mit \(\Sigma \subset \Gamma\),
    \item Bandalphabet \(\Gamma\) (endliche Menge von Symbolen),
    \item Übergangsfunktion \(\delta : Q \times \Gamma \to Q \times \Gamma \times \{L, R\}\),
    \item Startzustand \(q_0 \in Q\),
    \item Leerzeichen \(\shortsqcup \in \Gamma \setminus \Sigma\),
    \item Menge akzeptierender Zustände \(F \subseteq Q\).
\end{itemize}

\subsection{Graphische Darstellung}
Der Übergang \(\delta(q_0, X) = (q_1, Y, D)\) wird dargestellt als\\
\begin{center}
    \begin{tikzpicture}[main/.style = {draw, circle}]
        \node[main] (q) {$q_0$};
        \node[main] (p) [right=1.5cm of q] {$q_1$};
        \draw[->] (q) to node[above] {\(X/Y, D\)} (p);
    \end{tikzpicture}
\end{center}
wobei \(D \in \{L, R\}\) die Bewegungsrichtung des Kopfes angibt.

\subsection{Konfiguration}
Eine Zeichenreihe
\[
    K = X_1 X_2 \ldots X_{i-1}\, q\, X_i X_{i+1} \ldots X_n
\]
heisst Konfiguration von \(M\), wobei \(q \in Q\) der aktuelle Zustand ist und der Kopf über \(X_i\) steht. \\
\textbf{Startkonfiguration:} \(K_0 = q_0 X_1 X_2 \ldots X_n\) (Kopf auf erstem Zeichen).

\subsection{Berechnungsschritte}
Für \(\delta(q, X_i) = (p, Y, R)\) (Bewegung nach rechts):
\[
    \ldots X_{i-1}\, q\, X_i \ldots \vdash \ldots X_{i-1}\, Y\, p\, X_{i+1} \ldots
\]
Für \(\delta(q, X_i) = (p, Y, L)\) (Bewegung nach links):
\[
    \ldots X_{i-1}\, q\, X_i \ldots \vdash \ldots X_{i-2}\, p\, X_{i-1}\, Y\, X_{i+1} \ldots
\]

\subsection{Berechnung}
Eine Berechnung auf Eingabe \(w \in \Sigma^*\) ist eine endliche Folge
\[
    K_0 \vdash K_1 \vdash \cdots \vdash K_n
\]
wobei \(K_0\) die Startkonfiguration ist und aus \(K_n\) keine Bewegung mehr möglich ist. Kurzschreibweise: \(K_0 \vdash^* K_n\).

\subsection{Sprache einer TM}
Die von \(M\) akzeptierte Sprache ist
\[
    L(M) = \{ w \mid q_0 w \vdash^* \alpha p \beta,\; p \in F \text{ und } \alpha, \beta \in \Gamma^* \}.
\]
Eine Sprache, die von einer TM akzeptiert wird, heisst \textbf{rekursiv aufzählbar}. Die Klasse der rekursiv aufzählbaren Sprachen umfasst u.\,a.\ alle kontextfreien Sprachen.

\subsection{Erweiterungen von TM}

\subsubsection{TM mit Speichern}
In der Zustandssteuerung werden zusätzlich endlich viele Datenzustände gespeichert.
Ist nur eine Vereinfachung: \(Q_{\text{neu}} = Q \times \Gamma^k\) für \(k\) Speicher.

\subsubsection{TM mit mehreren Spuren}
Das Band hat \(n\) Spuren; jede Zelle speichert ein \(n\)-Tupel. Simulation durch
\(\Gamma_{\text{neu}} = \Gamma^n\).

\subsubsection{TM mit mehreren Bändern}
Die TM besitzt \(k\) Bänder mit je einem unabhängigen Lese-/Schreibkopf.
Übergangsfunktion:
\[
    \delta : Q \times \Gamma^k \to Q \times \Gamma^k \times \{R, L, S\}^k
\]
\textbf{Satz:} Jede Sprache, die von einer mehrbändigen TM akzeptiert wird, wird auch von einer TM mit nur einem Band akzeptiert (Simulation mit \(2k\) Spuren).

\subsubsection{Nichtdeterministische TM (NTM)}
Eine NTM hat die Übergangsfunktion
\[
    \delta : Q \times \Gamma \to \mathcal{P}(Q \times \Gamma \times D).
\]
\textbf{Satz:} Jede Sprache, die von einer NTM akzeptiert wird, wird auch von einer deterministischen TM akzeptiert (Simulation über Breitensuche aller Konfigurationen).

\subsection{Beschränkungen von TM}

\subsubsection{TM mit semi-unendlichem Band}
Das Band ist nur in eine Richtung unendlich. \\
\textbf{Satz:} Jede TM kann durch eine TM mit semi-unendlichem Band und zwei Spuren simuliert werden.

\subsubsection{Maschine mit mehreren Stacks}
Eine \(k\)-Stack-Maschine ist ein DKA mit \(k\) Stacks.
\[
    \delta : Q \times \Gamma_1 \times \cdots \times \Gamma_k \to Q \times \Gamma_1^* \times \cdots \times \Gamma_k^*
\]
\textbf{Satz:} Jede TM kann durch eine 2-Stack-Maschine simuliert werden (Bewegung durch Verschieben der Inhalte zwischen den beiden Stacks).

\subsubsection{Zählermaschine}
Eine Zählermaschine (Counter Machine) mit \(k\) Zählern entspricht einer \(k\)-Stack-Maschine, bei der Stacks durch Zähler in \(\mathbb{N}\) ersetzt werden. Nur \(= 0\) / \(\neq 0\) unterscheidbar; Inkrement/Dekrement um 1.\\
\textbf{Satz:} Jede TM kann durch eine 2-Zähler-Maschine simuliert werden. (via: TM \(\to\) 2-Stack \(\to\) 3-Zähler \(\to\) 2-Zähler).\\

\textbf{Stack als Zahl (Basis \(k\)):} Der Zustand \(s=X_1X_2\ldots X_n\) eines Stacks mit \(k-1\) Stacksymbolen kann durch einen Zähler \(z\) kodiert werden.
\begin{align*}
    z &= \sum_{i=0}^{n-1} X_{n-i} \cdot k^{n-i-1} \\
    \text{pop}&: z_{\text{new}} = z \div k, \quad X = z \mod k \\
    \text{push}(X_i)&: z_{\text{new}} = z \cdot k + X_i
\end{align*}
Der Dritte Zähler \(k\) dient zur Berechnung.\\

\textbf{3 Zähler \(\to\) 2 Zähler:} Die drei Zähler mit den Werten \(i,j,k\) werden als \(z = 2^i \cdot 3^j \cdot 5^k\) in einem Zähler kodiert. Mit dem zweiten Zähler wird gerechnet.


\section{Universelle Turingmaschine}
Eine \textbf{universelle TM} (UTM) simuliert jede beliebige TM $M$, sofern deren Codierung $\text{Cod}_M$ zusammen mit der Eingabe $w$ übergeben wird.

\subsection{Kodierung einer TM}
Jede TM lässt sich als Binärzeichenreihe über $\{0, 1\}$ kodieren:
\begin{enumerate}
    \item \textbf{Zustände:} $q_1$ = start, $q_2$ = akzeptierend, $q_3, \ldots, q_n$ = weitere Zustände.
    \item \textbf{Bandsymbole:} $X_1 \mathrel{\widehat{=}} 0$,\; $X_2 \mathrel{\widehat{=}} 1$,\; $X_3 \mathrel{\widehat{=}} \shortsqcup$,\; $X_4, \ldots , X_n$ = weitere Symbole.
    \item \textbf{Richtungen:} $D_1 \mathrel{\widehat{=}} L$,\; $D_2 \mathrel{\widehat{=}} R$.
    \item \textbf{Transition:} $\delta(q_i, X_j) = (q_k, X_l, D_m)$ wird kodiert als
    \[
        0^i 1\, 0^j 1\, 0^k 1\, 0^l 1\, 0^m
    \]
    \item \textbf{Gesamte TM:} Alle Transitionen $C_i$ werden durch \texttt{11} getrennt:
    \[
        C_1\, 11\, C_2\, 11\, C_3\, 11 \cdots
    \]
\end{enumerate}
Durch Voranstellen einer \texttt{1} entsteht eine eindeutige Zahl (Gödelnummer).

\subsection{Definition}
Eine TM $M_1$ heisst \textbf{universelle Turingmaschine}, wenn sie die Eingabe $\text{Cod}_M 111 w$ verarbeitet und das Ergebnis von $M$ auf $w$ liefert (für alle $M$, $w$). Die Folge \texttt{111} trennt die TM-Kodierung von der Eingabe $w$.


\section{Berechenbarkeitstheorie}

\subsection{Church-Turing-These}
\textbf{Intuitiv berechenbar:} Funktionen, die durch ein mechanisches Verfahren berechnet werden können.\\
\textbf{Turing-berechenbar:} Funktionen, die von einer TM berechnet werden können.\\
\textbf{Church-Turing-These:} Die Klasse der Turing-berechenbaren Funktionen stimmt mit der Klasse der intuitiv berechenbaren Funktionen überein.\\
\textbf{Gandy's These M:} Alles, was mit einer endlichen Maschine berechnet werden kann, ist bereits Turing-berechenbar. 
\textit{Nicht formal beweisbar.}

\subsection{Turing-berechenbare Funktionen}
Eine TM $T = (Q, \Sigma, \Gamma, \ldots)$ als partielle Funktion $T : \Sigma^* \to \Gamma^*$:
\[
    T(w) = \begin{cases} u & \text{falls } T \text{ auf } w \text{ mit } u \text{ auf dem Band anhält} \\ \uparrow & \text{falls } T \text{ bei Input } w \text{ nicht anhält} \end{cases}
\]
$f : \Sigma^* \to \Gamma^*$ heisst \textbf{Turing-berechenbar}, wenn es eine TM $T$ gibt mit $T(w) = f(w)$ für alle $w$ wo $f$ definiert ist.\\

\subsection{LOOP-Programme}
\textbf{Syntax:}
\begin{itemize}
    \item \textbf{Zuweisungen:} \texttt{x = y + c} und \texttt{x = y - c} wobei $x, y$ Variablen und $c \in \mathbb{N}$ ist.
    \item \textbf{Sequenzen:} \texttt{P ; Q}
    \item \textbf{Schleifen:} \texttt{Loop x Do P End}
\end{itemize}
\textbf{Semantik:} $P_k(n_1, \ldots, n_k) =$ Wert von $x_0$ nach Ablauf, mit $n_i$ in $x_i$.\\
\textbf{Konventionen:} $x_0$ und Hilfsvariablen mit 0 initialisiert. Subtraktion $\geq 0$.

\subsubsection{Bedingte Ausführung (IF x1=0 THEN P)}
\begin{verbatim}
x2 = x2 + 1;
Loop x1 Do x2 = x2 - 1 End;
Loop x2 Do P End
\end{verbatim}

\subsubsection{Verzweigung (IF x1=0 THEN P ELSE Q)}
\begin{verbatim}
x2 = x2 + 1; x3 = x3 + 1;
Loop x1 Do x2 = x2 - 1; x4 = x3 + 0 End;
Loop x2 Do P End;
Loop x4 Do Q End
\end{verbatim}

\subsection{WHILE-Programme}
Erweitern LOOP um: \texttt{While x > 0 Do P End}\\
\begin{itemize}
    \item Jedes LOOP-Programm ist ein WHILE-Programm
    \item WHILE-Programme terminieren \textbf{nicht} immer
\end{itemize}

\subsection{GOTO-Programme}
\textbf{Syntax:}
\begin{itemize}
    \item \texttt{Mk: xi = xj + c} / \texttt{Mk: xi = xj - c}
    \item \texttt{Mk: If xi = c Then Goto Mr}
    \item \texttt{Mk: Goto Mr}
    \item \texttt{Mk: HALT}
    \item Sequenzen: \texttt{P ; Q} (keine gemeinsamen Marker)
\end{itemize}

\subsection{Primitiv rekursive Funktionen}
\subsubsection{Grundfunktionen}
\begin{itemize}
    \item \textbf{Konstante:} $c_k^n : \mathbb{N}^n \to \mathbb{N}$, $c_k^n(x_1, \ldots, x_n) = k$
    \item \textbf{Nachfolger:} $\eta : \mathbb{N} \to \mathbb{N}$, $\eta(x) = x + 1$
    \item \textbf{Projektion:} $\pi_k^n : \mathbb{N}^n \to \mathbb{N}$, $\pi_k^n(x_1, \ldots, x_n) = x_k$
\end{itemize}

\subsubsection{Abschluss durch Einsetzung}
Sind $f, g_1, \ldots, g_k$ primitiv rekursiv mit $g_i : \mathbb{N}^n \to \mathbb{N}$, $f : \mathbb{N}^k \to \mathbb{N}$, dann auch:
\[
    h(\vec{x}) = f(g_1(\vec{x}), \ldots, g_k(\vec{x}))
\]

\subsubsection{Abschluss durch primitive Rekursion}
Sind $h : \mathbb{N}^n \to \mathbb{N}$ und $g : \mathbb{N}^{n+2} \to \mathbb{N}$ primitiv rekursiv, dann auch $f : \mathbb{N}^{n+1} \to \mathbb{N}$ mit:
\begin{align*}
    f(0, \vec{x}) &= h(\vec{x}) \\
    f(k+1, \vec{x}) &= g(f(k, \vec{x}), k, \vec{x})
\end{align*}

\subsubsection{Beispiele}
\textbf{Addition:}
$\text{Add}(0,y) = \pi_1^1(y)$, \quad $\text{Add}(x+1,y) = \eta(\pi_1^3(\text{Add}(x,y),x,y))$\\
\textbf{Multiplikation:}
$\text{Mult}(0,y) = c_0^1(y)$, \quad $\text{Mult}(x+1,y) = \text{Add}(\text{Mult}(x,y), \pi_2^2(x,y))$

\subsection{Unvollständigkeit von LOOP / Prim. Rek.}
LOOP-Programme und primitiv rekursive Funktionen sind \textbf{nicht} Turing-vollständig. Es gibt totale, berechenbare Funktionen, die nicht LOOP-berechenbar sind:

\subsubsection{Ackermannfunktion}
\begin{align*}
    a(0, m) &= m + 1 \\
    a(n+1, 0) &= a(n, 1) \\
    a(n+1, m+1) &= a(n, a(n+1, m))
\end{align*}
\textbf{Eigenschaften:}
\begin{itemize}
    \item Turing-berechenbar (WHILE-Programm möglich)
    \item \textbf{Nicht} primitiv rekursiv / LOOP-berechenbar
    \item Total (überall definiert) -- Beweis via Induktion über lexikograph.\ Ordnung auf $\mathbb{N} \times \mathbb{N}$
    \item Wächst schneller als jede primitiv rekursive Funktion
\end{itemize}

\subsubsection{LOOP-Interpreter}
$I : \mathbb{N}^2 \to \mathbb{N}$ mit $I(\langle P \rangle, x) = P_1(x)$ für jedes LOOP-Programm $P$.\\
\textbf{Satz:}
\begin{itemize}
    \item $I$ ist total und eindeutig (LOOP terminiert immer)
    \item $I$ ist Turing-berechenbar
    \item $I$ ist \textbf{nicht} LOOP-berechenbar
\end{itemize}
\textbf{Beweis (Widerspruch):} Angenommen $P_2$ ist LOOP-Interpreter. Definiere LOOP-Programm $Q$ mit $Q_1(x) = P_2(x,x) + 1$. Dann:
\[
    Q_1(\langle Q \rangle) = P_2(\langle Q \rangle, \langle Q \rangle) + 1 = Q_1(\langle Q \rangle) + 1 \quad \lightning
\]

\subsection{Überblick Berechnungsmodelle}
\[
    \underbrace{\text{LOOP} = \text{Prim.\ rek.}}_{\text{nicht Turing-vollständig}} \subsetneq \underbrace{\text{TM} = \text{WHILE} = \text{GOTO}}_{\text{Turing-vollständig}}
\]

\section{Entscheidbarkeit}

\subsection{Entscheidbarkeit}
Eine Sprache $A \subset \Sigma^*$ heisst \textbf{entscheidbar}, wenn eine TM $T$ existiert, die das Entscheidungsproblem $(\Sigma, A)$ löst:
\begin{itemize}
    \item $x \in A \Rightarrow T$ hält mit Bandinhalt ``1'' (Ja)
    \item $x \in \Sigma^* \setminus A \Rightarrow T$ hält mit Bandinhalt ``0'' (Nein)
\end{itemize}
$T$ muss bei \textbf{jeder} Eingabe $x \in \Sigma^*$ nach endlich vielen Schritten halten.\\
Äquivalent: $A$ ist entscheidbar $\Leftrightarrow$ $(\Sigma, A)$ mit einem WHILE-Programm lösbar (Entscheidungsverfahren).

\subsection{Semi-Entscheidbarkeit}
$A \subset \Sigma^*$ heisst \textbf{semi-entscheidbar}, wenn eine TM $T$ existiert mit:
\begin{itemize}
    \item $x \in A \Rightarrow T$ hält mit Bandinhalt ``1'' (Ja)
    \item $x \in \Sigma^* \setminus A \Rightarrow T$ hält \textbf{nie} an
\end{itemize}
Liefert nur positive Antworten; statt ``Nein'' keine Antwort.

\subsection{Konvention (natürliche Zahlen)}
$X \subset \mathbb{N}$ ist (semi-) entscheidbar $\Leftrightarrow$ $\{bin(x) \mid x \in X\} \subset \{0,1\}^*$ ist (semi-) entscheidbar.

\subsection{Zusammenhänge}
\begin{itemize}
    \item Jede entscheidbare Sprache ist auch semi-entscheidbar.
    \item $A$ ist entscheidbar $\Leftrightarrow$ sowohl $A$ als auch $\overline{A}$ sind semi-entscheidbar.
\end{itemize}
\textbf{Beweis ($\Leftarrow$):} Simuliere semi-Entscheidungsverfahren für $A$ und $\overline{A}$ schrittweise parallel:
\begin{verbatim}
INPUT(w); n = 0;
WHILE true DO
  n = n + 1;
  IF A(w,n) THEN return 1;
  IF B(w,n) THEN return 0
END
\end{verbatim}
$A(w,n)$: semi-EV von $A$ terminiert nach $n$ Schritten;\\$B(w,n)$: analog für $\overline{A}$.

\subsection{Abschlusseigenschaften}
\begin{itemize}
    \item $A$ entscheidbar $\Rightarrow$ $\overline{A}$ entscheidbar
    \item $A, B$ (semi-) entscheidbar $\Rightarrow$ $A \cup B$, $A \cap B$ (semi-) entscheidbar
\end{itemize}

\subsection{Charakterisierungen (Semi-Entscheidbarkeit)}
Für $A \subset \Sigma^*$ sind äquivalent:
\begin{itemize}
    \item $A$ ist rekursiv aufzählbar
    \item $A$ ist semi-entscheidbar
    \item $A$ ist Wertebereich einer totalen berechenbaren Funktion
    \item $A$ ist Definitionsbereich einer berechenbaren Funktion
\end{itemize}

\subsection{Reduktion}
$A \subset \Sigma^*$ heisst auf $B \subset \Gamma^*$ \textbf{reduzierbar} ($A \preceq B$), wenn es eine totale, Turing-berechenbare Funktion $F : \Sigma^* \to \Gamma^*$ gibt mit:
\[
    w \in A \Longleftrightarrow F(w) \in B
\]

\textbf{Transitivität:} $A \preceq B$ und $B \preceq C \Rightarrow A \preceq C$\\
\textbf{Satz:} Ist $B$ (semi-) entscheidbar und $A \preceq B$, dann ist auch $A$ (semi-) entscheidbar.\\
\textbf{Kontraposition:} Ist $A$ unentscheidbar und $A \preceq B$, dann ist auch $B$ unentscheidbar.\\
\textbf{Beweis (Entscheidbarkeit):}
\begin{verbatim}
INPUT(w); u = F(w); RETURN P(u)
\end{verbatim}

\subsection{Halteprobleme}

\subsubsection{Spezielles Halteproblem}
$H_S := \{w \in \{0,1\}^* \mid T_w \text{ angesetzt auf } w \text{ hält}\}$\\
\textbf{Gegeben:} Code $w$ einer TM $T_w$. \\
\textbf{Gefragt:} Hält $T_w$ auf eigenem Code $w$?

\subsubsection{Allgemeines Halteproblem}
$H := \{w\#x \in \{0,1,\#\}^* \mid T_w \text{ angesetzt auf } x \text{ hält}\}$\\
\textbf{Gegeben:} Code $w$ einer TM $T_w$ und Input $x$. \\
\textbf{Gefragt:} Hält $T_w$ auf Eingabewort $x$?

\subsubsection{Leeres Halteproblem}
$H_0 := \{w \in \{0,1\}^* \mid T_w \text{ angesetzt auf $\varepsilon$ Band hält}\}$\\
\textbf{Gegeben:} Code $w$ einer TM $T_w$. \\
\textbf{Gefragt:} Hält $T_w$ auf leerem Band?

\subsection{Unentscheidbarkeit der Halteprobleme}

\subsubsection{Unentscheidbarkeit von $H_S$}
Annahme: TM $T$ entscheidet $H_S$. Konstruiere TM $P$:
\begin{center}
\begin{tikzpicture}[node distance=2cm, auto]
  \node (T) [draw, rectangle] {TM $T$};
  \node (check) [draw, rectangle, right of=T] {Bandinhalt?};
  \node (set) [draw, rectangle, right of=check] {Band:=1};
  \node (halt) [draw, rectangle, right of=set] {Halt};
  
  \draw [->] (T) -- (check);
  \draw [->] (check) -- node[above] {0} (set);
  \draw [->] (set) -- (halt);
  \draw [->] (check) to [out=250,in=290,looseness=4] node[below] {1} (check);
\end{tikzpicture}
\end{center}
Sei $w$ der Code von $P$ ($T_w = P$):
\begin{itemize}
    \item $P$ auf $w$ hält $\Leftrightarrow T(w)=0 \Leftrightarrow T_w$ auf $w$ hält nicht $\Leftrightarrow P$ auf $w$ hält nicht $\lightning$
\end{itemize}

\subsubsection{Unentscheidbarkeit von $H$}
Reduktion $H_S \preceq H$ via $F(x) = x\#x$.\\
$w \in H_S \Leftrightarrow T_w$ hält auf $w \Leftrightarrow w\#w \in H \Leftrightarrow F(w) \in H$

\subsubsection{Unentscheidbarkeit von $H_0$}
Reduktion $H \preceq H_0$: Für Input $w\#x$ konstruiere TM $T'$, die zuerst $x$ auf das leere Band schreibt und sich dann wie $T_w$ verhält.\\
$T_w$ hält auf $x \Leftrightarrow T'$ hält auf leerem Band.

\subsection{Semi-Entscheidbarkeit der Halteprobleme}
$H_0$, $H_S$ und $H$ sind semi-entscheidbar.\\
Wegen $H_S \preceq H \preceq H_0$ genügt es, $H_0$ als semi-entscheidbar nachzuweisen: Universelle TM $U$ simuliert $T_w$ auf leerem Band; falls $U$ hält $\to$ Band $:= 1$.

\subsection{Satz von Rice}
Ist $R$ die Menge aller berechenbaren Funktionen und $S \subset R$ eine \textbf{echte, nichtleere} Teilmenge, dann ist
\[
    C(S) = \{w \in \{0,1\}^* \mid F_w \in S\}
\]
\textbf{unentscheidbar}.\\
\textbf{Konsequenzen:} Es ist im Allgemeinen unmöglich mechanisch zu prüfen, ob ein Programm\ldots
\begin{itemize}
    \item eine bestimmte Spezifikation erfüllt
    \item frei von Bugs ist
    \item bei jeder Eingabe terminiert
    \item dieselbe Funktionalität wie ein anderes Programm hat
\end{itemize}
\section{Komplexitätstheorie}

\subsection{Zeitkomplexität (TM)}
$\text{Time}_M(w) =$ Anzahl Konfigurationsübergänge von $M$ auf $w$.\\
\textbf{Worst-Case:} $\text{Time}_M(n) = \max\{\text{Time}_M(w) \mid |w| = n\}$

\subsection{$\mathcal{O}$-Notation (Landau-Symbole)}
\begin{itemize}
    \item $f \in \mathcal{O}(g)$: $\exists\, n_0, c$: $\forall\, n \geq n_0$: $f(n) \leq c \cdot g(n)$ \quad (obere Schranke)
    \item $f \in \Omega(g)$: $\exists\, n_0, d$: $\forall\, n \geq n_0$: $f(n) \geq \frac{1}{d} \cdot g(n)$ \quad (untere Schranke)
    \item $f \in \Theta(g)$: $f \in \mathcal{O}(g) \cap \Omega(g)$ \quad (asymptotisch gleich)
\end{itemize}
\textbf{Regeln:} $c \cdot f(n) \in \mathcal{O}(f(n))$;\; $c \in \mathcal{O}(1)$;\; $a_k n^k + \ldots + a_0 \in \mathcal{O}(n^k)$;\; transitiv. \\
\textbf{Abfolge}: $1 \ll \log n \ll n \ll n \log n \ll n^2 \ll 2^n \ll n!$

\subsection{Komplexitätsklassen}
\textbf{P:} Alle Probleme, die eine \textbf{det.\ TM} in Polynomzeit löst.\\
\textbf{NP:} Alle Probleme, die eine \textbf{NTM} in Polynomzeit löst. Äquivalent dazu: Menge aller Sprachen mit Polynomzeit-Verifizierer.\\
$\mathbf{P} \mathrel{\widehat{=}}$ Lösung \textit{finden} in Polyzeit;\;
$\mathbf{NP} \mathrel{\widehat{=}}$ Lösung \textit{verifizieren} in Polyzeit.\\
\textbf{Offene Frage:} $\mathbf{P} = \mathbf{NP}$? \quad Vermutung: $\mathbf{P} \neq \mathbf{NP}$.\\[3pt]
\textbf{NP-schwer:} $L$ ist NP-schwer, falls $\forall\, L' \in \mathbf{NP}$: $L' \preceq_p L$. Es existiert also eine polynomielle Reduktion von jedem NP-Problem auf $L$.\\
\textbf{NP-vollständig:} $L \in \mathbf{NP}$ \textit{und} $L$ ist NP-schwer.\\
Falls \textit{ein} NP-vollständiges $L \in \mathbf{P}$ $\Rightarrow$ $\mathbf{P} = \mathbf{NP}$.

\subsection{Polynomzeit-Verifizierer}
Ein Problem $A \in NP$ falls ein Verifizierungsproblem $A' \in P$ existiert, so dass 
$$ w \in A \Leftrightarrow \exists z: (|z| \leq p(|w|) \land z\#w \in A')$$

\subsection{Polynomielle Reduktion}
$A \preceq_p B$: Es existiert eine in Polynomzeit berechenbare Funktion $f$ mit $w \in A \Leftrightarrow f(w) \in B$.\\

\textbf{Bedeutet:} $B$ ist mind.\ so schwer wie $A$.\\
\textbf{Satz:} $P_1$ NP-schwer, $P_2 \in \mathbf{NP}$, $P_1 \preceq_p P_2$ $\Rightarrow$ $P_2$ NP-vollständig.



\subsection{SAT und Satz von Cook}
\textbf{SAT:} Ist eine KNF-Formel erfüllbar?\\
\textbf{Satz von Cook: SAT ist NP-vollständig.}\\
\textit{SAT $\in$ NP:} Zeuge $= $ Belegung; Verifikation in Polyzeit.\\
\textit{NP-schwer:} Berechnung einer polyzeit-NTM als KNF kodierbar.

\subsection{CLIQUE ist NP-vollständig}
\textbf{Clique:} $S \subseteq V$ paarweise verbunden. \\
\textbf{Problem:} $\exists$ Clique der Grösse $k$?\\
\textit{CLIQUE $\in$ NP:} Zeuge $= S$; Verifikation in $\mathcal{O}(n^2)$.\\
\textit{SAT $\preceq_p$ CLIQUE:} Aus KNF $\varphi$ mit $m$ Klauseln: $k = m$, Knoten $[i,j]$ pro Literal, Kanten zwischen versch.\ Klauseln falls Literale nicht komplementär.\\
$\varphi$ erfüllbar $\Leftrightarrow$ $G$ hat Clique der Grösse $k$.


\section{Chomsky-Hierarchie}
\begin{center}
     \begin{tikzpicture}[
          font=\scriptsize,
          box/.style={draw, rounded corners=2pt, inner sep=2pt, align=center},
          lbl/.style={anchor=north, align=center}
     ]
          % Typ 0: rekursiv aufzählbar (Turingmaschine)
          \node[box, fill=red!12, minimum width=4.1cm, minimum height=3cm] (t0) {};
          \node[lbl] at (t0.north) {\textbf{Typ 0} -- rekursiv aufzählbar\\(Turingmaschine)};

          % Typ 2: kontextfrei (Kellerautomat)
          \node[box, fill=green!22, minimum width=3.8cm, minimum height=2cm, anchor=south] (t2) at ([yshift=2pt]t0.south) {};
          \node[lbl] at (t2.north) {\textbf{Typ 2} -- kontextfrei\\(Kellerautomat)};

          % Typ 3: regulär (endlicher Automat)
          \node[box, fill=blue!18, minimum width=3.5cm, minimum height=1cm, anchor=south] (t3) at ([yshift=2pt]t2.south) {\textbf{Typ 3} -- regulär\\(endlicher Automat)};
     \end{tikzpicture}
\end{center}


\vfill

\rule{0.3\linewidth}{0.25pt}
\scriptsize

By Justin Iven Müller (2026)
\href{https://github.com/JustinIven/zhaw-cheatsheets}{github.com/JustinIven/zhaw-cheatsheets}


\end{multicols*}
\end{document}
